\documentclass[a4paper,11pt]{article}
\usepackage[utf8]{inputenc}
\usepackage[T1]{fontenc}
\usepackage[french]{babel}
\usepackage[left=2cm,right=2cm,top=2cm,bottom=2cm]{geometry}
\usepackage{xcolor}
\usepackage{hyperref}
\usepackage{fontawesome5}
\usepackage{parskip}

% --- Couleurs Sanofi ---
\definecolor{primary}{RGB}{102, 45, 145}
\hypersetup{colorlinks=true, urlcolor=primary}

\begin{document}

% --- EXPÉDITEUR ---
\noindent
\begin{minipage}[t]{0.5\textwidth}
    \textbf{Loïc Aron MBASSI EWOLO}\\
    Élève-Ingénieur Bac+5 -- Double diplôme ENSEA\\[3pt]
    \faMapMarker\ Cergy, France\\
    \faPhone\ (+33) 7 59 52 33 98\\
    \faEnvelope\ \href{mailto:loicaron.mbassiewolo@ensea.fr}{loicaron.mbassiewolo@ensea.fr}
\end{minipage}
\hfill
\begin{minipage}[t]{0.4\textwidth}
    \raggedleft
    Cergy, le 5 mars 2026
\end{minipage}

\vspace{1.2cm}

\hfill
\begin{minipage}[t]{0.5\textwidth}
    \raggedleft
    \textbf{Sanofi}\\
    Service Recrutement\\
    Gentilly / Lyon, France
\end{minipage}

\vspace{1cm}

\noindent\textbf{Objet :} Candidature -- Alternance Data Analyse, disponible dès septembre 2026

\vspace{0.8cm}

Madame, Monsieur,

Nettoyer un jeu de données médicales bruité, entraîner un modèle de classification, déployer une API Flask et construire un tableau de bord Streamlit pour que des non-techniciens puissent interpréter les résultats : c'est ce que j'ai réalisé lors d'IndabaX Africa 2025, compétition continentale d'IA centrée sur la santé, où mon équipe a remporté la première place. Ce projet illustre la façon dont je conçois la data analyse dans le secteur pharmaceutique : des données brutes aux décisions actionnables, avec le souci constant de la lisibilité et de la robustesse.

Élève-ingénieur en double diplôme ENSEA -- École Polytechnique de Yaoundé (Mathématiques Appliquées, Génie Logiciel), je suis actuellement en stage de recherche à l'INRIA Saclay (Équipe TRIBE, campus Institut Polytechnique de Paris), où je développe des pipelines Python d'analyse automatisée à grande échelle et des modèles NLP (BERT/Transformers) pour l'évaluation de la transparence de données textuelles. Ce travail sur la qualification et l'interprétabilité des données est transposable aux problématiques de Sanofi, où la fiabilité des référentiels de données conditionne directement la qualité des décisions.

Mon projet UROP 2024, mené avec des mentors Google DeepMind, portait sur la classification automatique d'anomalies cardiaques à partir d'images d'ECG : extraction de features signaux, traitement de données médicales hétérogènes, pipeline TensorFlow/Keras complet. En parallèle, mon travail de recherche au MILA (Institut Québécois d'IA, été 2025) sur l'interprétabilité des réseaux de neurones (phénomène de grokking, PyTorch) m'a formé à la rigueur mathématique nécessaire pour comprendre et améliorer les modèles prédictifs, au-delà de la simple mise en oeuvre.

Rejoindre Sanofi en alternance représente pour moi l'opportunité d'appliquer ces compétences à une échelle industrielle et à des données à fort enjeu clinique. Je suis disponible dès septembre 2026 et serais heureux d'échanger avec vous lors d'un premier entretien.

Je vous prie d'agréer, Madame, Monsieur, l'expression de mes salutations distinguées.

\vspace{0.8cm}

\textbf{Loïc Aron MBASSI EWOLO}\\
\textit{Élève-Ingénieur ENSEA / Polytechnique Yaoundé}

\end{document}
